\subsection{Experimental Design}
\label{sec:experimental-design}

To test the three formulated research hypotheses ($H_{1,1}$, $H_{1,2}$, and $H_{1,3}$), we adopt a $2\times2$ Factorial Design that crosses two independent variables: \textbf{Architecture Type} (\textit{Vanilla Kubernetes} vs. \textit{PANOPTES}) and \textbf{System State} (\textit{Idle Baseline} vs. \textit{Chaos/Fault Injection}). 

This factorial structure establishes four mutually exclusive experimental scenarios, detailed in Table \ref{tab:scenarios}. To achieve statistical power, mitigate the influence of transient virtualization noise, and fulfill the distributional assumptions required for robust inferential statistics, each scenario is executed across $N = 30$ independent runs. The sample size ($N = 30$) satisfies the Central Limit Theorem and provides sufficient power for both parametric tests and non-parametric rank-sum evaluations.

\input{chapters/05.methodology/tables/scenarios}

The technical configuration, analytical objective, and specific hypothesis addressed by each experimental scenario are defined as follows:

\begin{itemize}
    \item \textbf{Scenario 1 (Vanilla / Idle --- Control Baseline):} Standard, minimal Kubernetes environment running the target microservices workload without external observability agents. Telemetry persistence and monitoring rely exclusively on the native Kubernetes Metrics Server and container runtime streams. This scenario serves as the baseline reference point for resource footprint ($R_{\text{workload}}$), completely free from observer-induced overhead.
    
    \item \textbf{Scenario 2 (PANOPTES / Idle --- Static Overhead Evaluation):} Kubernetes environment running the exact same microservices workload alongside the instantiated PANOPTES telemetry stack (OpenTelemetry Collector, Prometheus, Grafana Loki, and Grafana Tempo). Comparing Scenario 2 directly against Scenario 1 isolates the net resource consumption of the observability platform ($R_{\text{PANOPTES}}$), providing the empirical data required to test the \textbf{Bounded Resource Overhead Hypothesis ($H_{1,3}$)} against the operational boundary ($\theta = 15\%$).
    
    \item \textbf{Scenario 3 (Vanilla / Stress --- Native Diagnostic Control):} The vanilla Kubernetes cluster is subjected to non-trivial fault injection campaigns orchestrated via Chaos Mesh (e.g., service-level network latency injection and latent cascade timeouts). Diagnostic investigations are conducted using native, unstructured CLI inspection workflows (\texttt{kubectl describe}, \texttt{kubectl logs}, \texttt{kubectl get events}). This scenario measures baseline diagnostic latency ($MTTDiag_{\text{Vanilla}}$) and native diagnostic accuracy ($Acc_{\text{Vanilla}}$).
    
    \item \textbf{Scenario 4 (PANOPTES / Stress --- Correlated Diagnostic Evaluation):} The PANOPTES-instrumented cluster is subjected to identical fault injection profiles under deterministic workload conditions. The operator conducts root-cause analysis leveraging the unified \textit{Single Pane of Glass}, utilizing metric-to-trace exemplars and Trace-to-Logs contextual deep linking. Evaluating Scenario 4 against Scenario 3 provides the comparative dataset used to test the \textbf{Diagnostic Efficiency Hypothesis ($H_{1,1}$)} and the \textbf{Diagnostic Accuracy Hypothesis ($H_{1,2}$)}.
\end{itemize}

By crossing system states with architecture types, this experimental design isolates confounding environmental variables, ensuring that measured differences in diagnostic speed, diagnostic success rate, and resource utilization are strictly attributable to the architectural contributions of PANOPTES.